%\begin{appendices}

\appendix
\setcounter{section}{0}% Reset numbering for sections
\renewcommand{\thesection}{\Alph{section}}% Adjust section printing (from here onward)

% ──────────────────────────────────────────────────────────────
\section{Árbol de Problemas}
% ──────────────────────────────────────────────────────────────

\label{anexo1}

\begin{figure}[h]
	\begin{center}
		\includegraphics[width=1.05\textwidth]{anexos/arbol_problemas.png}
		\caption[Árbol de problemas de la investigación]{Árbol de problemas identificados en la investigación sobre análisis de percepción nacional en contextos electorales mediante inteligencia electoral.\\
		Fuente: Elaboración propia.}
		\label{fig:arbol_problemas}
	\end{center}
\end{figure}

\clearpage

% ──────────────────────────────────────────────────────────────
\section{Árbol de Objetivos}
% ──────────────────────────────────────────────────────────────

\label{anexo2}

\begin{figure}[h]
	\begin{center}
		\includegraphics[width=1.05\textwidth]{anexos/arbol_objetivos.png}
		\caption[Árbol de objetivos de la investigación]{Árbol de objetivos que estructuran la investigación sobre análisis de percepciones con LLMs y detección dinámica de temas en ecosistemas digitales.\\
		Fuente: Elaboración propia.}
		\label{fig:arbol_objetivos}
	\end{center}
\end{figure}

\clearpage

% ──────────────────────────────────────────────────────────────
\begin{landscape}
    \section{Matriz de Consistencia}
    \label{anexo3}
    
    \begingroup
    \renewcommand\arraystretch{1.3}
    \small
    \begin{longtable}{p{2.8cm} p{3.2cm} p{3.2cm} p{2.8cm} p{2.8cm} p{2.8cm} p{2.8cm}}
        
        \caption[Matriz de consistencia de la investigación]{Matriz de consistencia que articula problema general, objetivo general, hipótesis general, variables, dimensiones e indicadores de la investigación.}\\
        \label{tbl:matriz_consistencia} \\
        
        \specialrule{.1em}{.05em}{.05em}
        \centering\textbf{TÍTULO} & \centering\textbf{PROBLEMA} & \centering\textbf{OBJETIVO} & \centering\textbf{HIPÓTESIS} & \centering\textbf{VARIABLE} & \centering\textbf{DIMENSIÓN} & \centering\arraybackslash\textbf{INDICADOR} \\
        \specialrule{.1em}{.05em}{.05em}
        \endfirsthead

        \caption[]{Matriz de consistencia (continuación).}\\
        \specialrule{.1em}{.05em}{.05em}
        \centering\textbf{TÍTULO} & \centering\textbf{PROBLEMA} & \centering\textbf{OBJETIVO} & \centering\textbf{HIPÓTESIS} & \centering\textbf{VARIABLE} & \centering\textbf{DIMENSIÓN} & \centering\arraybackslash\textbf{INDICADOR} \\
        \specialrule{.1em}{.05em}{.05em}
        \endhead
        
        % Variable Independiente - Usamos {=} para que respete el ancho y no se sobreescriba
        \multirow{8}{=}{Inteligencia Electoral con LLM: Minería de Sentimientos y Detección Dinámica de Temas Relevantes para analizar la Percepción Nacional en Ecosistemas Digitales Electorales}
        & \multirow{5}{=}{PG: ¿Es posible analizar la percepción nacional mediante un sistema de Inteligencia Electoral basado en LLMs y ecosistemas multi-fuente?}
        & \multirow{5}{=}{OG: Analizar la percepción nacional mediante un sistema basado en LLMs, minería de sentimientos y detección de temas.}
        & \multirow{5}{=}{HG: El uso de LLMs permite analizar con mayor precisión la percepción nacional en entornos multi-fuente.}
        & \multirow{5}{=}{Independiente: Sistema de Inteligencia Electoral con LLM}
        & \multirow{5}{=}{Arquitectura del pipeline}
        & \multirow{5}{=}{Integración de módulos ABSA, BERTopic, Neo4j y GraphRAG} \\
        & & & & & & \\
        & & & & & & \\
        & & & & & & \\
        \hline
        
        % Variable Dependiente
        & PE1: ¿Qué enfoques existen para el análisis de opinión pública mediante IA?
        & OE1: Analizar enfoques existentes en la literatura.
        & HE1: La revisión de literatura permite identificar metodologías aplicables.
        & \multirow{10}{=}{Dependiente: Percepción pública nacional}
        & Sentimiento ciudadano
        & Score continuo $[-1, 1]$ por candidato \\
        \cline{2-4} \cline{6-7}
        
        & PE2: ¿Cómo integrar múltiples fuentes de datos para representar la percepción?
        & OE2: Diseñar una arquitectura multi-fuente.
        & HE1: La integración de múltiples fuentes mejora representatividad.
        &
        & Metainformación del corpus
        & Volumen por fuente y período \\
        \cline{2-4} \cline{6-7}
        
        & PE3: ¿Cuál es el impacto de LLMs en análisis de sentimientos y detección de temas?
        & OE3: Evaluar desempeño de LLMs en análisis.
        & HE2: Los LLMs mejoran precisión del análisis.
        &
        & Temas relevantes
        & Coherencia semántica $C_v$ y cobertura \\
        \cline{2-4} \cline{6-7}
        
        & PE4: ¿Cómo modelar relaciones discurso mediático-percepción ciudadana?
        & OE4: Implementar sistema basado en grafos de conocimiento.
        & HE3: GraphRAG captura relaciones complejas.
        &
        & Evolución temporal
        & Indicadores de evolución temporal \\
        
        \specialrule{.1em}{.05em}{.05em}
    \end{longtable}
    \endgroup
\end{landscape}

\clearpage

% ──────────────────────────────────────────────────────────────
\section{Configuración de Parámetros del Módulo ABSA con LLM}
% ──────────────────────────────────────────────────────────────

\label{anexo4}

\begin{table}[h]
	\centering	
	\small
	\begin{tabular}{p{4cm} p{5cm} p{4cm}}
		\specialrule{.1em}{.05em}{.05em}
		\textbf{Parámetro} & \textbf{Valor/Configuración} & \textbf{Justificación} \\
		\specialrule{.1em}{.05em}{.05em}
		
		Modelo LLM & Gemini Flash 1.5 & Balance entre velocidad de inferencia y calidad de embeddings contextuales \\
		\hline
		
		Rol del Prompt & Analista político experto & Orienta el modelo hacia el dominio electoral latinoamericano \\
		\hline
		
		Ejemplos \textit{Few-Shot} & 4 ejemplos & De acuerdo con \cite{pr_gujral2024llm_elections}, mínimo 3 ejemplos mejora precisión en contextos no anglosajones \\
		\hline
		
		Escala de Sentimiento & Continua $[-1, 1]$ & Captura intensidad del sentimiento más allá de polaridad binaria \\
		\hline
		
		Umbral Clasificación & Positivo: $S > 0.2$, Negativo: $S < -0.2$, Neutro: $[-0.2, 0.2]$ & Reduce ruido de sentimientos débiles \\
		\hline
		
		Estrategia \textit{Early Stopping} & 3 iteraciones sin mejora F1 & Evita sobreajuste y desperdicio de tokens API \\
		\hline
		
		Formato de Respuesta & JSON con validación esquema & Estructuración consistente de candidatos, sentimientos, temas y relaciones \\
		\hline
		
		Temperatura/Creatividad & Baja (0.2-0.3) & Reproducibilidad y consistencia en clasificación \\
		
		\specialrule{.1em}{.05em}{.05em}
	\end{tabular}
	\caption[Configuración de parámetros del módulo ABSA]{Parámetros finales configurados para el módulo de análisis de sentimientos a nivel de aspecto (ABSA) usando Gemini Flash 1.5.}
	\label{tbl:config_absa}
\end{table}

\clearpage

% ──────────────────────────────────────────────────────────────
\section{Configuración de Parámetros del Módulo BERTopic}
% ──────────────────────────────────────────────────────────────

\label{anexo5}

\begin{table}[h]
	\centering
	\small
	\begin{tabular}{p{4cm} p{5cm} p{4cm}}
		\specialrule{.1em}{.05em}{.05em}
		\textbf{Parámetro} & \textbf{Valor} & \textbf{Justificación} \\
		\specialrule{.1em}{.05em}{.05em}
		
		Modelo Embedding & \texttt{paraphrase-multilingual-MiniLM-L12-v2} & Mejor balance entre calidad semántica y velocidad para español \cite{mc_reimers2019sbert} \\
		\hline
		
		Dimensión Embedding & 384 & Dimensión estándar del modelo multilingüe seleccionado \\
		\hline
		
		Reducción Dimensionalidad & UMAP a 5 dimensiones & Mejora visualización y clustering sin perder información semántica \\
		\hline
		
		Algoritmo Clustering & HDBSCAN inductivo & No requiere especificar $k$ a priori; excluye automáticamente ruido \\
		\hline
		
		min\_cluster\_size & 5 & Previene fragmentación extrema en clusters políticos \\
		\hline
		
		min\_samples & 5 & Densidad mínima para considerar punto como núcleo \\
		\hline
		
		Etiquetado de Temas & Gemini + c-TF-IDF & LLM genera etiquetas semánticamente ricas de 15 términos top \\
		\hline
		
		Coherencia Semántica & $C_v$ & Métrica de validación para evaluar calidad temática \\
		\hline
		
		\textit{Early Stopping} & Sin mejora $C_v$ & Detiene optimización de HDBSCAN al platear coherencia \\
		
		\specialrule{.1em}{.05em}{.05em}
	\end{tabular}
	\caption[Configuración de parámetros de BERTopic]{Parámetros de configuración del módulo BERTopic para detección dinámica de temas electorales basado en publicaciones en redes sociales.}
	\label{tbl:config_bertopic}
\end{table}

\clearpage

% ──────────────────────────────────────────────────────────────
\section{Esquema del Grafo de Conocimiento Electoral en Neo4j}
% ──────────────────────────────────────────────────────────────

\label{anexo6}

\begingroup
\renewcommand\arraystretch{1.3}
\small
\setlength{\LTpost}{0pt}

\begin{longtable}{p{2.5cm} p{3.5cm} p{6.5cm}}
    \caption[Esquema del grafo de conocimiento electoral]{Especificación completa de nodos, propiedades y relaciones implementadas en el grafo de conocimiento electoral en Neo4j AuraDB.} \label{tbl:esquema_grafo} \\
    
    \specialrule{.1em}{.05em}{.05em}
    \centering\textbf{TIPO NODO} & \centering\textbf{PROPIEDADES} & \centering\arraybackslash\textbf{DESCRIPCIÓN} \\
    \specialrule{.1em}{.05em}{.05em}
    \endfirsthead

    \specialrule{.1em}{.05em}{.05em}
    \centering\textbf{TIPO NODO} & \centering\textbf{PROPIEDADES} & \centering\arraybackslash\textbf{DESCRIPCIÓN} \\
    \specialrule{.1em}{.05em}{.05em}
    \endhead
    
    Candidato & id, nombre, partido, color\_hexa, estatus & Candidatos presidenciales en el proceso electoral \\ \hline
    Tema & id, nombre, etiqueta\_llm, coherencia\_cv, palabras\_clave, frecuencia & Temas emergentes detectados por BERTopic \\ \hline
    Fuente & id, tipo, nombre, url, alcance & Portales de noticias, redes sociales, canales de video \\ \hline
    Evento & id, fecha, tipo, descripción, impacto & Debates, mítines, declaraciones críticas \\ \hline
    Período & id, fecha\_inicio, fecha\_fin, nombre & Ventanas temporales de análisis (semanal, mensual) \\ \hline
    
    \specialrule{.1em}{.05em}{.05em}
    \textbf{TIPO RELACIÓN} & \multicolumn{2}{p{10cm}}{\centering\textbf{METADATOS Y DESCRIPCIÓN}} \\
    \specialrule{.1em}{.05em}{.05em}
    
    MENCIONA & timestamp, frecuencia, sentimiento\_promedio & Documento menciona a candidato \\ \hline
    EXPRESA\_SENTIMIENTO & score\_continuo, fecha, fuente\_tipo & Candidato-Tema con polaridad \\ \hline
    ASOCIADO\_A & confianza, coocurrencia & Tema asociado a evento electoral \\ \hline
    PUBLICADO\_EN & fecha, alcance, interacciones & Documento publicado en fuente específica \\ \hline
    PERTENECE\_A & período & Evento/documento asignado a período temporal \\ \hline
    RELACIONADO\_CON & similitud\_temática & Temas relacionados \\ \hline
    REACCIONA\_A & timestamp, sentimiento & Ciudadano reacciona a evento \\ \hline
    INFLUYE & intensidad & Evento influye en percepción de candidato \\ 
    \specialrule{.1em}{.05em}{.05em}
\end{longtable}
\endgroup

\clearpage

% ──────────────────────────────────────────────────────────────
\section{Diccionario de Palabras Clave Electorales}
% ──────────────────────────────────────────────────────────────

\label{anexo7}

\begin{table}[h]
	\centering
	\small
	\setlength{\LTpost}{0pt}
	\begin{longtable}{p{3cm} p{7cm} p{3cm}}
		\specialrule{.1em}{.05em}{.05em}
		\centering\textbf{CATEGORÍA} & \centering\textbf{TÉRMINOS} & \centering\arraybackslash\textbf{CANTIDAD} \\
		\specialrule{.1em}{.05em}{.05em}
		\endhead
		
		Procesos Electorales & elección, voto, votación, candidato, campaña, debate, mitin, encuesta, sondeo & 9 \\
		\hline
		
		Términos Políticos & gobierno, política, legislatura, congreso, república, democracia, parlamento & 7 \\
		\hline
		
		Instituciones & ONPE, JNE, RENIEC, INE, INEI, MINDEF & 6 \\
		\hline
		
		Temas de Campaña & economía, seguridad, educación, salud, corrupción, gobernabilidad, infraestructura & 7 \\
		\hline
		
		Hashtags Electorales & \#ElecciónesPerú2024, \#Voto2024, \#DebatePresidencial, \#CuéntaTuVoto & 4+ \\
		\hline
		
		Candidatos & Nombres completos, apodos, siglas de partidos & Dinámico según proceso \\
		
		\specialrule{.1em}{.05em}{.05em}
		
	\end{longtable}
	\caption[Diccionario de palabras clave electorales]{Vocabulario controlado utilizado para filtrar publicaciones relevantes en el corpus multi-fuente durante el período de análisis electoral.}
	\label{tbl:diccionario_palabras_clave}
\end{table}

\clearpage

% ──────────────────────────────────────────────────────────────
\section{Métodos Estadísticos para Validación de Modelos}
% ──────────────────────────────────────────────────────────────

\label{anexo8}

\subsection{Métricas de Desempeño para ABSA}

Para la evaluación del módulo de análisis de sentimientos a nivel de aspecto, se emplearon las siguientes métricas estadísticas en escala $[0, 1]$:

\begin{equation}
\text{Exactitud} = \frac{\text{TP} + \text{TN}}{\text{TP} + \text{TN} + \text{FP} + \text{FN}}
\end{equation}

\begin{equation}
\text{Precisión} = \frac{\text{TP}}{\text{TP} + \text{FP}}
\end{equation}

\begin{equation}
\text{Cobertura (Recall)} = \frac{\text{TP}}{\text{TP} + \text{FN}}
\end{equation}

\begin{equation}
\text{F1-Score Macro} = \frac{1}{n}\sum_{i=1}^{n}\frac{2 \cdot \text{Precisión}_i \cdot \text{Cobertura}_i}{\text{Precisión}_i + \text{Cobertura}_i}
\end{equation}

Donde TP (Verdaderos Positivos), TN (Verdaderos Negativos), FP (Falsos Positivos), FN (Falsos Negativos) representan los conteos de clasificación por clase.

\subsection{Métricas de Coherencia Temática para BERTopic}

Para la validación de la detección dinámica de temas, se utilizó la métrica $C_v$ de coherencia semántica:

\begin{equation}
C_v = \frac{1}{|\text{Temas}|}\sum_{t}\text{cohere}_V(W_t)
\end{equation}

Donde $W_t$ es el conjunto de palabras clave del tema $t$, y $\text{cohere}_V$ es la función de coherencia vectorial. Rangos de interpretación:
\begin{itemize}
	\item $C_v < 0.4$: Coherencia baja
	\item $0.4 \leq C_v < 0.6$: Coherencia media
	\item $C_v \geq 0.6$: Coherencia alta
\end{itemize}

\clearpage

% ──────────────────────────────────────────────────────────────
\section{Timeline Experimental de CRISP-DM}
% ──────────────────────────────────────────────────────────────

\label{anexo9}

\begin{table}[h]
	\centering
	\small
	\begin{tabular}{p{2.5cm} p{3.5cm} p{2.5cm} p{4cm}}
		\specialrule{.1em}{.05em}{.05em}
		\centering\textbf{FASE} & \centering\textbf{ACTIVIDADES CLAVE} & \centering\textbf{DURACIÓN} & \centering\arraybackslash\textbf{ENTREGABLES} \\
		\specialrule{.1em}{.05em}{.05em}
		
		1. Comprensión del Negocio & Definición de problema, objetivos e hipótesis; Revisión de literatura & 2 semanas & Documento de especificaciones \\
		\hline
		
		2. Comprensión de Datos & Construcción de corpus (noticias, X, YouTube); Análisis exploratorio & 3 semanas & Datasets en capas Bronze y Silver \\
		\hline
		
		3. Preparación de Datos & Limpieza, normalización, tokenización; Segmentación por período & 2 semanas & Corpus preprocesado validado \\
		\hline
		
		4. Modelamiento & Config. ABSA, BERTopic, Neo4j, GraphRAG; Calibración de parámetros & 4 semanas & Modelos entrenados y validados \\
		\hline
		
		5. Evaluación & Métricas de desempeño, análisis comparativo, validación cualitativa & 2 semanas & Reporte de resultados \\
		\hline
		
		6. Despliegue & Prototipo funcional con datos en tiempo real; Documentación técnica & 1 semana & Sistema operacional \\
		
		\specialrule{.1em}{.05em}{.05em}
	\end{tabular}
	\caption[Timeline experimental de las fases CRISP-DM]{Cronograma de implementación de las seis fases de la metodología CRISP-DM aplicada en la investigación.}
	\label{tbl:timeline_crisp_dm}
\end{table}


